% ============================================================
% 分片：s03d  对应原文 PDF 第 45--46 页
% 负责章节：§3.2.2 Example: Landau--Ginzburg $B$-model（例：Landau--Ginzburg $B$ 模型）
% 译者：translator-07   审校状态：初稿
% ============================================================
% [跨页接续: 原文第 44 页末句"…We will assume $A$ is $\mathbb{Z}$-graded by the fermion
%  number operator $F$"已由 s03c 译出；本分片自该句的"so that we can write …"起接续，
%  按 Lead 的边界协调不重复"fermion number operator $F$"。]

% ==== 原文 PDF 第 45 页 ====

使得我们可以写 $A = \bigoplus_{j\in\mathbb{Z}} A^j$，
其中 $A^j$ 是 $A$ 中满足 $[F,A^j] = jA^j$ 的元素。我们只关心两个超荷中的一个 $Q$：
它与哈密顿量对易，并把 $A^j$ 经 $[Q,-]$ 送到 $A^{j+1}$，因此称 $Q$ 为次数 $1$ 的导子。
事实上，它在算子乘积上的作用恰是一个导子：
\begin{equation}\tag{3.25}
[Q,ab] = [Q,a]\,b + (-1)^{F(a)} a\,[Q,b] .
\end{equation}
% [Tutor Note: 式 (3.25) 是分次 Leibniz 律。$Q$ 是奇算子（$F(Q)=1$），当它「穿过」$a$ 时须付出
%  Koszul 符号 $(-1)^{F(a)}$；把 $Q$ 写成左乘对易子 $[Q,-]$ 之后，(3.25) 正是「$Q$ 是 $A$ 的
%  DGA 微分」的等价表述。易错点：$F(a)$ 只取 $a$ 的宇称（$\mathbb{Z}$ 次数 mod 2），
%  不要误写成 $F(b)$ 或 $F(ab)$。]
本讨论中始终用 $[-, -]$ 表示分次对易子（即依算子的宇称，它或为反对易子、或为对易子）。
相对于这个指定的超荷 $Q$，$A$ 成为以 $d_A = [Q,-]$ 为微分的 DGA。
类似地，态空间（Hilbert 空间）作为该代数带相容分次与微分的表示，是一个 DG 模。
满足 $[Q,a]=0$ 的 $a\in A$ 常称为（$\frac{1}{2}$-）BPS 算子，
相应的上同调类记作 $[a]\in H^\bullet(A)$。

DGA 的上同调上的高阶运算，可以由上同调类的提升之乘积（及更高结合子）与这些乘积（及更高结合子）
的提升之间的同伦来构造。此处涉及许多选择，但所得的 $A_\infty$ 代数彼此等价，即拟同构。
首先选取映射 $f_1:H^\bullet(A)\to A$，把每个上同调类提升为表示它的上闭链。
于是对任意 $[a],[b]$，差 $f_1([a][b]) - f_1([a])f_1([b])$ 按构造在上同调中平凡
（即该差为 $Q$-恰当的），我们选取一个见证这一 $Q$-恰当性的映射
$f_2:H^\bullet(A)^{\otimes 2}\to A$：
% [Tutor Note: 推导补足：(3.26) 的含义是，$f_1([a][b])$ 与 $f_1([a])f_1([b])$ 都是 $Q$-闭的，
%  二者之差因而 $Q$-闭；它在 $H^\bullet(A)$ 中为零，故必可写成 $Q$ 作用在某个上链上的像，$f_2$ 就是
%  这个「原像」。这也说明 $f_2$ 并不典范：作替换 $f_2\to f_2+Qg$ 只改变高阶运算的代表元，
%  不改变上同调类。]
\begin{equation}\tag{3.26}
Q f_2([a],[b]) = f_1([a][b]) - f_1([a]) f_1([b]) .
\end{equation}
元素 $[a]$、$[b]$、$[c]$ 的三元积 $\mu_3([a],[b],[c])$ 则满足
\begin{equation}\tag{3.27}
\begin{aligned}
f_1(\mu_3([a],[b],[c])) &= f_2([a],[b])f_1([c]) - (-1)^{F([a])} f_1([a])f_2([b],[c]) \\
&\quad + f_2([a][b],[c]) - f_2([a],[b][c]) + Q(\dots) ,
\end{aligned}
\end{equation}
% [Tutor Note: 推导补足：(3.27) 是 $n=3$ 的 $A_\infty$ 关系 (3.24) 在「取 $H^\bullet(A)$ 上同调类」
%  之后的形式。右端各项来自把结合子 $([a][b])[c]-[a]([b][c])$ 的提升张开：
%  $f_2([a],[b])f_1([c])$ 与 $f_1([a])f_2([b],[c])$ 分别修正两种结合顺序，
%  $f_2([a][b],[c]) - f_2([a],[b][c])$ 修正乘积本身的提升，末尾的 $Q(\dots)$ 项是不影响
%  上同调类的 $Q$-恰当部分。$(-1)^{F([a])}$ 来自把 $f_1([a])$ 移过 $f_2$ 时的 Koszul 符号。]
即 $\mu_3([a],[b],[c])$ 是右端表达式的上同调类。

这一过程继续下去。给定到 $i=n-2$ 为止的提升 $f_i$ 与到 $i=n-1$ 为止的乘积 $\mu_i$，
可选取一个 $(n-1)$ 元提升 $f_{n-1}:H^\bullet(A)^{\otimes n-1}\to A$，它见证某个表达式的 $Q$-恰当性；
该表达式包含 $f_1(\mu_{n-1}([a],[b],\dots))$ 与一个涉及低阶 $i$ 提升与乘积的（显式）表达式之差。
$n$ 元积 $\mu_n([a],[b],\dots)$ 于是取作某个（显式）表达式的上同调类，
该表达式涉及 $f_1,f_2,\dots,f_{n-1}$ 与低阶 $i$ 乘积。更多细节见 [77]。
诸映射 $f_i$ 定义了一个从 $\{H^\bullet(A),\mu_n\}$ 到 $(A,Q)$ 的 $A_\infty$ 态射，
它在上同调层面是恒等映射。这同样是拟同构。

把目前为止所见总结如下：$\mathcal{N}=2$ 量子力学中的局域算子，在选定扭变超荷 $Q$ 之后，
自然带有 DGA 结构。过渡到 $Q$-上同调类（或过渡到扭变理论），意味着丢弃物理理论的大量信息。
然而，这些丢失的信息中有相当一部分仍可编码为上同调类上的更高 $n$ 元运算：
这些高阶运算强到足以把物理的局域算子恢复到 $A_\infty$ 等价为止。

\subsubsection{例：Landau--Ginzburg $B$ 模型}

下面在略更有趣的情境中显式计算上述高阶乘积的一个例子：二维理论一维边界上的局域算子代数。
这同时也说明
% ==== 原文 PDF 第 46 页 ====
% [跨页接续: 原文第 45 页末"This will also illustrate"与第 46 页首
%  "how $A_\infty$ algebras can transfer across dualities via quasi-isomorphisms"为同一句，
%  此处合并译出。]
$A_\infty$ 代数如何通过拟同构在对偶之间转移。
我们考虑三个手征超场 $\Phi^n = \phi^n + \dots$（$n=1,2,3$）的二维 Landau--Ginzburg 模型，
超势为 $W = \Phi^1\Phi^2\Phi^3$，然后取 $B$ 型扭变。
与 $B$ 型扭变相容的边界条件范畴（即 $B$ 膜范畴）被 Kontsevich 提出为超势 $W$ 的（导出）矩阵分解范畴
[78, Section 7]。附录 C 简要回顾了矩阵分解如何作为 Landau--Ginzburg 模型中 $B$ 膜范畴的描述出现。
在本例中，该（导出）范畴由三个对象生成：矩阵分解 $E = \phi^1$、$J = 2\partial_{\phi^1}W$ 及其置换。
特别地，任何矩阵分解 $(E^a, J_a)$ 都可实现为这三个矩阵分解的直和，或可能是其形变与/或直和项。

% [背景补充: 「矩阵分解」指满足 $EJ = JE = W\cdot\mathrm{id}$ 的一对映射 $(E,J)$，即把超势「开平方」；
%  Kontsevich 据此把 Landau--Ginzburg 模型的 $B$ 膜实现为 $W$ 的矩阵分解的（导出）范畴，
%  对象之间的态射复形即该范畴的 Hom 复形。]

我们考虑矩阵分解 $E=(\phi^1,\phi^2,\phi^3)$ 与 $J=\frac{2}{3}(\phi^2\phi^3,\phi^1\phi^3,\phi^1\phi^2)$；
这一选择与体手征多重态上的一个 Dirichlet 边界条件\footnote{%
这类 Dirichlet 矩阵分解由 Dyckerhoff 详细研究过（他称之为 \emph{stabilized residue fields}），
并且已知当 $W$ 具有孤立奇点时，它们本身就生成矩阵分解的导出范畴 [79, Theorem 4.1]。
此外，其自同态 DGA $A$ 与弯曲结合代数 $R_W = (\mathbb{C}[\phi^a|], W)$ Koszul 对偶
[80, Proposition 3.2]。遗憾的是，此处的 $W$ 奇性太强，这些结果并不适用——超曲面 $W = 0$
沿除子 $\phi^1 = \phi^2 = 0$、$\phi^1 = \phi^3 = 0$ 与 $\phi^2 = \phi^3 = 0$ 中的每一条都是奇异的。%
}相关——如此选取的 $E$ 意在强制 $\phi^n = 0$。具体地，我们发现该 Dirichlet 边界条件上的局域算子
DGA 为
\begin{equation}\tag{3.28}
A = \mathbb{C}[\phi^n, \bar\gamma_n, \partial_{\bar\gamma_n}], \qquad
Q = \phi^n\bar\gamma_n + \frac{2}{3}W_n\partial_{\bar\gamma_n},
\end{equation}
% [校订: 原文第 46 页 (3.28) 两处生成元均印作有横线的 $\bar\gamma_n$，与 (3.29) 一致；
%  初稿曾照文字层误录为 $\gamma_n$，现按页面图订正。]
% [存疑: 脚注中弯曲结合代数印作 $R_W = (\mathbb{C}[\phi^a|], W)$，方括号内的竖线疑为原稿排版冗余；
%  Dyckerhoff 的规范记号通常为 $(\mathbb{C}[\phi^a], W)$。此处照抄原样。]
其中 $W_n := \partial_{\phi^n}W$。值得注意的是，该矩阵分解被置换指标 $n$ 的 $S_3$ 对称性所保持。
一个直接的运算表明，三个费米子算子
\begin{equation}\tag{3.29}
\psi_1 = \bar\gamma_1 - \frac{1}{3}\left(\phi^2\partial_{\bar\gamma_3} + \phi^3\partial_{\bar\gamma_2}\right) \qquad
\psi_2 = \bar\gamma_2 - \frac{1}{3}\left(\phi^1\partial_{\bar\gamma_3} + \phi^3\partial_{\bar\gamma_1}\right) \qquad
\psi_3 = \bar\gamma_3 - \frac{1}{3}\left(\phi^2\partial_{\bar\gamma_3} + \phi^3\partial_{\bar\gamma_2}\right)
\end{equation}
% [原文疑误: (3.29) 中 $\psi_3$ 括号内的两项印作 $\phi^2\partial_{\bar\gamma_3} + \phi^3\partial_{\bar\gamma_2}$，
%  与 $\psi_1$ 括号内的表达式完全相同。按 $S_3$ 置换对称性应为
%  $\phi^1\partial_{\bar\gamma_2} + \phi^2\partial_{\bar\gamma_1}$。判断依据：把
%  $\{\bar\gamma_m,\partial_{\bar\gamma_n}\} = \delta_{mn}$ 代入，用印出的 $\psi_3$ 算得 $\{\psi_1,\psi_3\}$
%  除 $Q$-恰当项（如 $2\bar\gamma_1\bar\gamma_3$）外的 $c$ 数部分为 $-\frac{1}{3}\phi^2$，
%  与 (3.30) 给出的 $-\frac{2}{3}\phi^2$ 相差因子 $2$；按对称性更正后二者一致。
%  此处照抄原样，建议按对称性更正。]
是 $Q$-闭的但不是 $Q$-恰当的。此外，它们生成 $Q$-上同调。非零的反对易子为
\begin{equation}\tag{3.30}
\{\psi_1,\psi_2\} = -\frac{2}{3}\phi^3 \qquad
\{\psi_1,\psi_3\} = -\frac{2}{3}\phi^2 \qquad
\{\psi_2,\psi_3\} = -\frac{2}{3}\phi^1 ,
\end{equation}
% [Tutor Note: 计算 (3.30) 需用 $\{\bar\gamma_m,\partial_{\bar\gamma_n}\} = \delta_{mn}$
%  （把 $\partial_{\bar\gamma_n}$ 视为与 $\bar\gamma_n$ 共轭的费米子自由度），并注意 $\psi_n$ 为费米子
%  算子，故 $\psi_1\psi_2 = -\psi_2\psi_1$；由此 $\{\psi_1,\psi_2\}$ 才能算出非零的 $c$ 数结果。
%  易错点：若漏掉 (3.29) 括号内两项的符号或指标配对，会得到 $\{\psi_m,\psi_n\}=0$ 的错误结论。]
它们全部是 $Q$-恰当的：$[\psi_n][\psi_m] = -[\psi_m][\psi_n]$。特别地，诸 $[\psi_n]$ 生成一个外代数：
$H^\bullet(A)\cong\wedge^\bullet\mathbb{C}^3$。尽管如此，下面将看到存在非平凡的三元积
$\mu_3([\psi_1],[\psi_2],[\psi_3])$。
% [Tutor Note: $H^\bullet(A)\cong\wedge^\bullet\mathbb{C}^3$ 说明边界算子的上同调环在外代数层面
%  是「自由」的：通常的乘积结构几乎不含信息，全部内容都藏在更高阶运算中。这正是 (3.33) 中
%  $\mu_3\neq 0$ 的意义——若只看 DGA 的乘积，这个 $B$ 膜边界条件与三个自由费米子无从区分。]

为计算三元积 $\mu_3([\psi_1],[\psi_2],[\psi_3])$，需要构造提升 $f_1$、$f_2$。
首先选取如下上同调类的一元提升
\begin{equation}\tag{3.31}
\begin{aligned}
f_1(1) &= 1 \qquad
f_1([\psi_n]) = \psi_n \qquad
f_1([\psi_1][\psi_2]) = \psi_1\psi_2 + \frac{1}{3}\phi^3 \\
f_1([\psi_1][\psi_2][\psi_3]) &= \psi_1\psi_2\psi_3 + \frac{1}{3}\left(\phi^1\psi_1 - \phi^2\psi_2 + \phi^3\psi_3\right)
\end{aligned}
\end{equation}
其中其余的提升由 $S_3$ 作用得到\footnote{%
这些提升与 $H^\bullet(A)$ 上的代数关系相容，这一点很重要。特别地，我们不能随意地提升
$[\psi_1][\psi_2]$ 与 $[\psi_2][\psi_1]$：它们被要求满足
$f_1([\psi_1][\psi_2]) = -f_1([\psi_2][\psi_1])$。本质上，$f_1$ 的这一选取对应于提升到「正规序」上链，
例如 $f_1([\psi_1][\psi_2]) = \frac{1}{2}(\psi_1\psi_2 - \psi_2\psi_1) = \psi_1\psi_2 + \frac{1}{3}\phi^3$。%
}。为构造二元提升 $f_2$，需要选取 $f_1([a][b])$ 与 $f_1([a])f_1([b])$ 之间的一个同伦。
例如，我们发现
% [跨页接续: 原文第 46 页末句"For example, we find that"续至第 47 页开头
%  "$f_1([\psi_1][\psi_2]) - f_1([\psi_1])f_1([\psi_2]) = \frac{1}{3}\phi^3 = \frac{1}{3}Q\partial_{\gamma_3}$
%  implying that we can choose $\dots$"，由 s03e 接续。]
